\section{Free products and coherent amalgamation}\label{sec:app-products}
The common free-product theorem gives N005 Theorem~3.8, including both directions of flag transitivity and residual connectedness and the firmness equivalence. For Theorem~3.23 the amalgamation isomorphism must be one fixed map respecting every indexed common parabolic. The common tree theorem and its converse then prove all standard intersections, both directions of flag transitivity, the residue assertions and all panel conditions. This is the coherent interpretation explicitly used on source page~23. The distinctions below explain why the map and the chamber stabilizer cannot be suppressed.

The trivial-Borel condition cannot be omitted. Take both inputs to
be the rank-one geometry $(C_2^2,(C_2))$. Each has exactly two
vertices and is thin and flag-transitive. Their free-product geometry
has two types, and every panel is infinite by
the common free-product panel formula. This is a counterexample to the
unrestricted thinness assertion in Piedade--Tranchida,
arXiv:2505.24662v1, Corollary 3.9. Their definition permits
nonfaithful actions; imposing a regular action with trivial Borel
recovers the asserted thin case. the common free-product theorem also removes
the extra Borel assumption from their converse flag assertion.


\subsection{The small free products}
Types remain part of a coset system even if they contain one element.
Thus Example 3.10's rank-one factor $(1,(1))$ produces the original
parabolics together with one equal to the whole group. The result
is the cone on the input, with one new universal element, and is
not type-preservingly isomorphic to the input. A rank-zero factor
would give the asserted identity. In Example 3.11 the two factors
$(C_2,(1))$ give the Bass--Serre tree of $C_2*C_2$: its vertices
are cosets of the two factors, its edges are group elements, and
both vertex degrees are two. This is the infinite bivalent tree.

For an $n$-gon ($n\geq3$) and the rank-one $C_2$ factor, the
presentation is
$\langle a_1,a_2,a_3\mid a_i^2=1,(a_1a_2)^n=1\rangle$.
The subgroups obtained by deleting one generator are precisely
Example 3.12's maximal parabolics. Under the classical Coxeter
reflection-triangle realization it is the tessellation with angles
$\pi/n,0,0$. Amalgamating an $n$-gon and an $m$-gon over their
shared reflection instead gives
\[
 \langle a_1,a_2,b_3\mid a_1^2=a_2^2=b_3^2=1,
 (a_1a_2)^n=(a_2b_3)^m=1\rangle.
\]
Its maximal parabolics are the three deleted-generator subgroups,
as in Example 3.26. For $n,m\geq3$ the angles $\pi/n,\pi/m,0$
define a hyperbolic triangle. More generally this interpretation
requires $1/n+1/m<1$; the algebraic amalgam formula has no such
geometric inequality.

Example 3.25's empty shared type set means amalgamation over the
identified Borel subgroups, and is a free product if both are
trivial. Identifying the two entire systems gives that system.
The parenthetical assertion that thin flag-transitive systems always
have trivial Borel requires a faithful acting group. It is false
for arbitrary specified group actions; the common exact panel
formula keeps the Borel explicitly.

\subsection{The amalgamation map is part of compatibility}
The positive amalgamation theorem uses one chosen isomorphism of the
common parabolics carrying every indexed intersection to the matching
one. Separate abstract isomorphisms, as written in Definition 3.16,
do not imply this coherence for an arbitrary chosen amalgamation.
For example take two standard rank-two geometries of $C_2^2$, with
coordinate generators $a,b$ and corresponding maximal subgroups
$\langle a\rangle,\langle b\rangle$. Let all types be shared, so
the amalgamated subgroup is the whole group, but identify the two
copies by the automorphism interchanging $a,b$. Every pair appearing
in Definition 3.16 is abstractly isomorphic. Nevertheless the output
maximal subgroup of either type is
$\langle a,b\rangle=C_2^2$, since it combines different coordinate
lines under the chosen identification. The output consists of two
universal singleton types and is not firm, while both inputs are thin.

Thus Lemma 3.17's step from abstract isomorphism to equality of the
embedded intersections needs the specified coherent map. The explicit
assumption on page 23 that all the maximal parabolics are special
supplies exactly that condition. All positive uses here retain it;
none applies a gluing theorem to arbitrary abstractly isomorphic
parabolic data. This example qualifies the literal shorthand, without
refuting the coherent-map theorem proved in the common paper.

\subsection{The incompatible symmetric-group parabolics}
In Example 3.18 the literal generators $a_0,a_1,a_2$ satisfy the
Coxeter $[4,3]$ relations and $(a_0a_1a_2)^3=1$. In the signed
permutation model of $[4,3]$, of order 48, this cube of a Coxeter
element is central inversion. Hence $A_3=\langle a_0,a_1,a_2\rangle$
has order at most 24. The assignment
\[
 a_0=(12)(34),\quad a_1=(23),\quad a_2=(34),\quad a_3=(14)(23)
\]
satisfies every relator of the full source group in $S_4$. The first
three images generate $S_4$, proving $A_3\cong S_4$ exactly.
The pair $a_0,a_1$ generates $D_8$: its defining order bound is
attained in this quotient. Since $a_3$ commutes with both,
$A_2=\langle a_0,a_1,a_3\rangle$ has order dividing 16.
Lagrange's theorem now bounds $|A_2\cap A_3|$ by
$\gcd(16,24)=8$, and this intersection contains the displayed
$D_8$. Thus equality holds. On the other side the Coxeter $[3,3,4]$
parabolic $B_4$ is $[3,3]\cong S_4$ and
$B_2\cap B_4=\langle b_0,b_1\rangle\cong D_6$.
These different intersection orders prove the claimed incompatibility
without retaining the classification of the whole regular polytope.


\subsection{Why local incidence information is necessary}


The following examples concern the unrestricted subgroup statements of
N005 Lemmas~3.15(a) and~3.29(a). They do not by themselves disprove
the later geometric theorems, whose parabolic hypotheses supply additional
structure. In particular, two cosets meeting the same third coset need not
meet each other inside that third coset. This is precisely the inference
that must be justified when comparing reduced expressions.

\begin{proposition}\label{X-prop:amalgam}
Special subgroups of a free product with amalgamation need not have the
intersection asserted in N005 Lemma~3.15(a).
\end{proposition}
\begin{proof}
Let
\[
 A=\langle a,c\mid a^2=c^2=[a,c]=1\rangle,\qquad
 B=\langle b,c\mid b^2=c^2=[b,c]=1\rangle,
 \qquad C=\langle c\rangle.
\]
Both factors are Klein four groups and $G=A*_C B$ has presentation
\[
 \langle a,b,c\mid a^2=b^2=c^2=[a,c]=[b,c]=1\rangle.
\]
Set $D_A=\langle a\rangle$, $D_B=\langle b\rangle$,
$E_A=\langle ac\rangle$ and $E_B=\langle bc\rangle$.
Each of these four groups intersects $C$ trivially. Thus
$D=\langle a,b\rangle$ and $E=\langle ac,bc\rangle$ are special
subgroups in the exact sense of that lemma, with trivial amalgamated
subgroups. Yet
\[
 ab=(ac)(bc)\in D\cap E.
\]
The word $ab$ is nontrivial by the reduced word theorem, whereas
$D_A\cap E_A=D_B\cap E_B=1$. The formula in the lemma would therefore
give $D\cap E=1$, a contradiction.

For completeness, the intersection can be determined exactly. The displayed
presentation identifies $G$ with $(C_2*C_2)\times C_2$. In this product,
$D$ is the first factor and $E$ is the graph of the homomorphism
$C_2*C_2\longrightarrow C_2$ sending both distinguished involutions to
the nonidentity element. Its kernel is the infinite cyclic subgroup
$\langle ab\rangle$, by alternating reduced words. Hence
$D\cap E=\langle ab\rangle$.
\end{proof}

\begin{proposition}\label{X-prop:hnn}
Special subgroups of an HNN extension need not have the intersection
asserted in N005 Lemma~3.29(a).
\end{proposition}
\begin{proof}
Use the same Klein four group $A=\langle a,c\rangle$ and identify
$C=\langle c\rangle$ with itself by the identity. Its HNN extension is
\[
 G=\langle a,c,t\mid a^2=c^2=[a,c]=[c,t]=1\rangle.
\]
The groups $D=\langle a,t\rangle$ and $E=\langle ac,t\rangle$ are
special: both base subgroups meet the associated subgroup $C$ trivially.
The asserted intersection formula would give $D\cap E=\langle t\rangle$.
However,
\[
 w=ata=(ac)t(ac)\in D\cap E.
\]
To prove $w\notin\langle t\rangle$, map $G$ to $S_3\times C_2$ by
$a\mapsto((12),1)$, $t\mapsto((23),1)$ and
$c\mapsto(1,\epsilon)$, where $\epsilon$ generates $C_2$.
Every defining relation is satisfied. The image of $w$ is $((13),1)$,
which is outside the cyclic subgroup generated by $((23),1)$.
\end{proof}

The same finite group separates the nontrivial word in the first example:
send the generators $a,b,c$ to the respective permutations $(12),(23),(45)$ in $S_5$. The factor groups remain
embedded, their intersection is $\langle(45)\rangle$, and $ab$ has order
three. These explicit homomorphisms are elementary finite witnesses for the stated nonmembership. No finite quotient is used to assert an equality of subgroups in the infinite groups. The earlier optional GAP verification remains provenance; it is not a proof premise.

The relevant positive property is stronger than arbitrary subgroup
compatibility. In a flag-transitive coset geometry, pairwise incident
cosets have a common representative. Thus a coset of a parabolic that
meets two other incident parabolic cosets meets their intersection.
The examples above lack this property for the three local subgroups
$C,\langle a\rangle,\langle ac\rangle$. A corrected gluing argument
must retain this local incidence information.
